\begin{proposition}[Hodge 规范下异常类的一点消失判据]\label{prop:conclusion-anomaly-hodge-onepoint-criterion}
\leanverified{paper\_conclusion\_anomaly\_hodge\_onepoint\_criterion}
在定义 \ref{def:pom-bounded-slice} 的任意有界切片 \(\mathbf{PW}^{\ast}_{\le}\) 上，设投影词 \(w\) 的异常类
\[
[{\rm Anom}(w)]\in H^2(K;A_G)
\]
落在 torsion-free 扇区。则由定理 \ref{thm:pom-pw-minnorm-counterterm-hodge} 可取其 Hodge 规范代表 \(\mathfrak h(w)\)，并且下列命题等价：
\begin{enumerate}
  \item \([{\rm Anom}(w)]=0\)；
  \item \(\mathfrak h(w)(x_0)=0\) 于某一认证点 \(x_0\)；
  \item \(w\) 的交叉异常坐标与 Euler 商不变量同时消失。
\end{enumerate}
因此，在 torsion-free 口径下，异常类的消失可由一点 Hodge 数据与有限个离散不变量完全判定。
\end{proposition}
\begin{proof}
定理 \ref{thm:pom-pw-minnorm-counterterm-hodge} 给出最小范数 strictification 势与 Hodge 规范代表；推论 \ref{cor:pom-anom-one-point-zero-test} 把零异常压缩为一点判别；而定理 \ref{thm:pom-euler-defect-quotient-coordinate} 说明 torsion-free 扇区中的异常类可由交叉异常坐标与 Euler 商不变量完全坐标化。三者合并即得结论。证毕。
\end{proof}

\begin{theorem}[strictifiability 的有限证书完备性]\label{thm:conclusion-strictifiability-finite-certificate-completeness}
\leanverified{paper\_conclusion\_strictifiability\_finite\_certificate\_completeness}
在有界切片 \(\mathbf{PW}^{\ast}_{\le}\) 上，strictifiability 不是无限维同伦问题，而是有限证书问题。更准确地，对每个投影词 \(w\) 可定义有限签名
\[
\Xi_m(w):=
\bigl(
G^{\mathrm{BC}}_m(w),\
\mathrm{Eul}(w),\
\mathrm{Cross}(w),\
\Sigma_m(w)
\bigr),
\]
其中 \(G^{\mathrm{BC}}_m(w)\) 为普适可交换化商坐标，\(\mathrm{Eul}(w)\) 为 Euler 商不变量，\(\mathrm{Cross}(w)\) 为交叉异常坐标，而 \(\Sigma_m(w)\) 为注 \ref{rem:pom-finite-signature-sigmam} 的有限语义签名。则
\[
w\ \text{strictifiable}
\quad\Longleftrightarrow\quad
\Xi_m(w)\ \text{为平凡证书}.
\]
换言之，strictifiability 在有界切片上已被压缩为一个有限可核验的判定问题。
\end{theorem}
\begin{proof}
由命题 \ref{prop:conclusion-anomaly-hodge-onepoint-criterion}，异常类在 torsion-free 口径下是否消失，可由 \(\mathrm{Eul}(w)\) 与 \(\mathrm{Cross}(w)\) 的有限数据决定。另一方面，定理 \ref{thm:pom-sem-conservative} 说明 \(\mathsf{Sem}\) 在有界切片上反射自然同构，而注 \ref{rem:pom-finite-signature-sigmam} 已把该语义输出压缩为有限签名 \(\Sigma_m(w)\)。再结合普适 strictification 商 \(G^{\mathrm{BC}}_m(w)\) 的有限坐标化，strictifiability 所需的异常残差、群商残差与语义残差都被 \(\Xi_m(w)\) 吸收；其平凡性恰等价于 strictification 的存在。证毕。
\end{proof}

\begin{corollary}[判等与优化的同一有限签名原则]\label{cor:conclusion-finite-signature-unifies-equivalence-and-optimization}
\leanverified{paper\_conclusion\_finite\_signature\_unifies\_equivalence\_and\_optimization}
在同一有界切片上，对任意两个投影词 \(w_1,w_2\)，若
\[
\Xi_m(w_1)=\Xi_m(w_2),
\]
则：
\begin{enumerate}
  \item \(w_1,w_2\) 属于同一 strictification 类；
  \item 对任意 admissible 代价函子
  \[
  \mathcal C:\mathbf{PW}\to\mathbf{Res},
  \]
  二者具有相同的 Pareto 前沿与同一规范代表函子输出。
\end{enumerate}
因此，在有界切片上，“判等”与“优化”并不需要两套证书，而由同一有限签名统一控制。
\end{corollary}
\begin{proof}
第一点是定理 \ref{thm:conclusion-strictifiability-finite-certificate-completeness} 的直接推论。第二点由定义 \ref{def:pom-cost-functor}、命题 \ref{prop:pom-cost-optimization-decidable} 与推论 \ref{cor:pom-cost-canonical-representative-functor} 给出：一旦 \(\Xi_m\) 相同，\(\mathsf{Sem}\) 输出与 strictification 数据都落在同一自然同构类上，而 admissible 代价函子与规范代表函子只依赖该类与其有限前沿，故两者的 Pareto 前沿与规范代表一致。证毕。
\end{proof}

\begin{theorem}[算子值完全正可见实现与 strictification 的解析等价]\label{thm:conclusion-operator-caratheodory-cp-visible-equivalence}
固定分辨率 \(m\)，令
\[
\mathcal A_m:=\CC[G_m^{\mathrm{BC}}]
\cong
\bigoplus_{\pi\in\widehat{\mathcal A_m}}M_{d_\pi}(\CC).
\]
设 \(F_m:\mathbb D\to\mathcal A_m\) 为处于零异常 strictification 口径下的可见解析对象，并记其不可约通道分量为 \(F_{m,\pi}\)。则下列条件等价：
\begin{enumerate}
  \item \(F_m\) 为 \(\mathcal A_m\)-值 Carath\'eodory 函数，且全部块 Toeplitz 截断满足
  \[
  \mathbf T_N^{\mathcal A_m}(F_m)\succeq 0
  \qquad(\forall N\ge 0);
  \]
  \item 对每个 \(\pi\in\widehat{\mathcal A_m}\)，分量 \(F_{m,\pi}\) 都是矩阵值 Carath\'eodory 函数；
  \item 对每个 \(\pi\)，其 Cayley--Schur 变换
  \[
  S_{m,\pi}(w):=\bigl(F_{m,\pi}(w)-I\bigr)\bigl(F_{m,\pi}(w)+I\bigr)^{-1}
  \]
  满足
  \[
  \mathrm{sq}_-\!\bigl(S_{m,\pi}\bigr)=0;
  \]
  \item \(F_m\) 诱导一个统一的完全正可见实现。
\end{enumerate}
故在 strictification 可见商上，完全正可见实现、块 Toeplitz--PSD、逐通道 Carath\'eodory 与零负平方指数是同一现象的四种表述。
\end{theorem}
\begin{proof}
定理 \ref{thm:xi-operator-herglotz-toeplitz-completeness} 与其逐块分解说明 \(\textup{(1)}\Leftrightarrow\textup{(2)}\)。矩阵值 Cayley--Schur 对应把通常 Carath\'eodory 类精确送到通常 Schur 类，因此 \(\textup{(2)}\Leftrightarrow\textup{(3)}\)。在零异常 strictification 口径下，定理 \ref{thm:xi-gbc-maximal-cp-visible-realization} 把 \(\mathcal A_m\)-值 Herglotz/Toeplitz 正性与统一完全正可见实现对齐，故 \(\textup{(1)}\Leftrightarrow\textup{(4)}\)。而定理 \ref{thm:xi-cp-visible-implies-kappa0} 正给出 \(\textup{(4)}\Rightarrow\textup{(3)}\) 的显式通道化版本。证毕。
\end{proof}

\begin{corollary}[完全正不可实现性的有限维否证器]\label{cor:conclusion-cp-visible-nonrealizability-has-finite-refuter}
\leanverified{paper\_conclusion\_cp\_visible\_nonrealizability\_has\_finite\_refuter}
在定理 \ref{thm:conclusion-operator-caratheodory-cp-visible-equivalence} 的设定下，若 \(F_m\) 不能诱导完全正可见实现，则必存在
\begin{enumerate}
  \item 一个不可约通道 \(\pi\in\widehat{\mathcal A_m}\)，以及
  \item 一个有限阶 \(N\ge 0\),
\end{enumerate}
使得对应的有限块 Toeplitz 条件失败：
\[
\mathbf T_N(F_{m,\pi})\not\succeq 0.
\]
因此，完全正不可实现性必有有限维否证器，而不存在“只在无限层才显影”的纯软失败。
\end{corollary}
\begin{proof}
这是定理 \ref{thm:conclusion-operator-caratheodory-cp-visible-equivalence} 的否命题。若不存在完全正可见实现，则至少有一个条件失败；由 \(\textup{(1)}\Leftrightarrow\textup{(2)}\) 的逐通道块分解与推论 \ref{cor:xi-abelian-channel-negativity-localization}，失败必可局域到某一不可约通道与某一有限阶 Toeplitz 主子式。故得到有限维否证器。证毕。
\end{proof}

\begin{theorem}[热带极限中的 Pareto 几何刚性]\label{thm:conclusion-tropical-pareto-rigidity-on-bounded-slice}
在有界切片 \(\mathbf{PW}^{\ast}_{\le}\) 上，设
\[
\mathcal C:\mathbf{PW}\to\mathbf{Res}
\]
为 admissible 代价函子，并对温度门施加正文的低温热带化口径。则：
\begin{enumerate}
  \item 复合代价严格满足热带乘法律
  \[
  \mathcal C(w_2\circ w_1)=\mathcal C(w_2)\otimes \mathcal C(w_1),
  \]
  即空间—时间成本按 min-plus 规则累积；
  \item 每个 strictification 类的代价像 \(\mathcal C([w])\subset\mathbb T^{(2)}\) 具有唯一的 Pareto 下确界；
  \item 对任何被规范代表函子选中的 Pareto 极小实现，不存在另一条同类 \(2\)-态射能同时严格降低两条成本坐标。
\end{enumerate}
因此，strictification 类内的 Pareto 几何在热带极限中是刚性的：代价可比较、可组合，但不存在“双坐标同时免费下降”的隐藏方向。
\end{theorem}
\begin{proof}
第一点只是定义 \ref{def:pom-cost-functor} 的复合公理与正文温度门热带化原则的结论层重述。第二点由命题 \ref{prop:pom-cost-optimization-decidable} 得到：在有界切片上，\(\mathcal C([w])\) 的帕累托前沿是有限集，从而其在 \(\mathbb T^{(2)}\) 的点态预序下有唯一坐标下确界。第三点由推论 \ref{cor:pom-cost-canonical-representative-functor} 的规范代表最优性推出：若存在同类 \(2\)-态射同时严格降低两坐标，则规范代表便不是帕累托极小元，矛盾。证毕。
\end{proof}

\begin{conjecture}[残差细分与定位增强的二律背反]\label{conj:conclusion-residue-refinement-localization-tradeoff}
设 \(P(q,u)\) 为正文 \(9.38\) 的联合压力面，\(\mathcal R\) 为任意非平凡 residue 细分。则对所有 \(q>1\) 与充分低温的 \(u\)，应同时满足：
\begin{enumerate}
  \item 碰撞压力严格下降：
  \[
  P_{\mathcal R}(q,u)<P(q,u);
  \]
  \item 至少一个可见通道的反向 KL 指纹严格增大；
  \item 因而不存在一种非平凡 residue 细分，能够同时实现“碰撞压缩严格改善”与“anomaly 局部化成本在全通道上不增”。
\end{enumerate}
若该猜想成立，则 strictification、compression 与 localization 之间将满足一个真正的 no-free-lunch 定理：打散碰撞的收益，必以某个通道的定位增强代价为代偿。
\end{conjecture}
